% ---------------------------------------------------------------------------
% LRP per-rule ablation -- CAMBRIA Capstone, extension experiment 2
% Build:  pdflatex lrp_rule_ablation.tex   (twice, for refs)
% ---------------------------------------------------------------------------
\documentclass[11pt,a4paper]{article}

\usepackage[margin=1in]{geometry}
\usepackage{amsmath,amssymb}
\usepackage{booktabs}
\usepackage{longtable}
\usepackage{xcolor}
\usepackage{graphicx}
\usepackage[colorlinks=true,linkcolor=blue!50!black,urlcolor=blue!50!black,
            citecolor=blue!50!black]{hyperref}
\usepackage{listings}
\usepackage{enumitem}

\lstset{basicstyle=\ttfamily\footnotesize, breaklines=true,
        frame=single, framesep=4pt, rulecolor=\color{black!25},
        backgroundcolor=\color{black!3}, columns=fullflexible}

\newcommand{\jlens}{J\nobreakdash-lens}
\newcommand{\rlens}{R\nobreakdash-lens}
\newcommand{\cfg}[1]{\texttt{#1}}

\title{\textbf{Which LRP Rule Carries the \rlens?}\\[0.3em]
       \large A per-rule ablation of the three stop-gradient rules\\
       that separate the \rlens{} from the \jlens}
\author{CAMBRIA Capstone --- Extension Experiment 2}
\date{August 2026}

\begin{document}
\maketitle

\begin{abstract}
\noindent
The \rlens{} improves on the \jlens{} by installing three layer-wise relevance
propagation (LRP) stop-gradient rules before fitting. The original work reports the
combined effect and never separates the rules, and its fitting code was never
released. We rebuilt the fitting harness, verified it reproduces the released
artifacts, and fitted a lens for each of the eight subsets of the three rules on two
models under an otherwise identical recipe.

\textbf{Result.} Only the \emph{half-rule} contributes. It recovers $53\%$ of the
full \rlens{} advantage on Qwen3.5\nobreakdash-27B and $112\%$ on
Qwen3.5\nobreakdash-4B. The \emph{LN-rule} alone scores \emph{below} the \jlens{}
baseline on both models. The \emph{identity-rule} has no reliable effect --- it
changes sign between models.

\textbf{Secondary result.} Weight-space similarity to the released \rlens{} fails to
predict lens quality. Two arms are mispredicted in \emph{sign}. We report this
separately because, unlike the attribution itself, it does not depend on corpus size.
\end{abstract}

\tableofcontents
\newpage

% ===========================================================================
\section{The problem}

\subsection{What the \rlens{} is}

The \jlens{} approximates the map from an intermediate residual stream state to the
model's final state by its average input--output Jacobian,
\begin{equation}
  J_\ell \;=\; \mathbb{E}_{x,t}\!\left[\frac{\partial h_{\text{final}}}{\partial h_\ell}\right],
  \qquad
  \operatorname{lens}(h) \;=\; \operatorname{unembed}(J_\ell\, h),
\end{equation}
with the expectation taken over prompts $x$ and token positions $t$ from a
pretraining-like corpus.

The \rlens{} changes nothing about this definition. It changes the \emph{gradients}
the expectation is taken over, by installing three LRP stop-gradient rules into the
backward pass before fitting. The forward pass is bit-identical; only $\partial
h_{\text{final}}/\partial h_\ell$ differs.

\subsection{The three rules}

\begin{longtable}{@{}p{0.16\textwidth}p{0.34\textwidth}p{0.42\textwidth}@{}}
\toprule
\textbf{Rule} & \textbf{What it changes} & \textbf{Motivation} \\
\midrule
\endhead
\cfg{ln\_rule} & Detach the RMSNorm normalizer, so gradient does not flow through the
  normalization scale factor. & Treats normalization as a fixed rescaling rather than
  a function of the input. \\
\addlinespace
\cfg{identity\_rule} & Replace the SiLU backward with $\sigma(x)$, i.e.\ treat
  $\mathrm{SiLU}(x)$ as locally $x \cdot \sigma(x)$ with $\sigma$ constant. & Removes
  the derivative-of-gate term, an LRP-style ``identity'' relevance assignment. \\
\addlinespace
\cfg{half\_rule} & Split the SwiGLU product gradient $50/50$ between its two factors
  ($\beta = 0.5$). & Distributes relevance evenly across a multiplicative
  interaction, rather than by local partial derivative. \\
\bottomrule
\end{longtable}

\subsection{The question, and why nobody could answer it}

The published work reports the \rlens{} as a package. Three rules go in, one improved
lens comes out, and the contribution of each is unknown. Answering this requires
fitting lenses with each rule subset --- which requires fitting code that was never
released. Only the final \emph{weights} were published.

This is the gap the experiment fills: eight lenses per model, one per subset of
$\{\text{ln}, \text{identity}, \text{half}\}$, fitted on identical prompts with an
identical recipe, so that any difference is attributable to the rules alone.

% ===========================================================================
\section{Prerequisite: rebuilding and verifying the fitting harness}

Every result below is downstream of one question: \emph{does our reimplementation of
the fitting procedure actually reproduce theirs?} If not, the sweep measures
properties of our bug rather than properties of the rules.

\subsection{Recovering the recipe}

The released artifacts carry provenance metadata, which we extracted rather than
guessed. All four released lenses (4B and 27B, \jlens{} and \rlens) record:

\begin{center}
\begin{tabular}{@{}ll@{}}
\toprule
\texttt{dataset\_id}   & \texttt{NeelNanda/pile-10k} \\
\texttt{n\_prompts}    & \textbf{25} \\
\texttt{docs\_consumed}& 25 \\
\texttt{t\_max}        & 128 tokens \\
\texttt{skip\_first}   & 4 positions \\
\texttt{target\_layer} & 30 (4B) / 62 (27B), i.e.\ penultimate \\
\texttt{weighting}     & uniform \\
\texttt{corpus\_mode}  & pretrain \\
\texttt{git\_commit}   & \texttt{"modal"} \\
\bottomrule
\end{tabular}
\end{center}

\noindent
We further inferred the prompt rows: \texttt{docs\_consumed == n\_prompts} implies no
documents were skipped, so the draw is pile-10k rows $[0,25)$.

We also verified that the released files append $J_{\text{target}} = I$ exactly
($\|J_{30} - I\|_F = 0$ in both released 4B lenses) and mirrored that.

\subsection{Bit-exactness of the patched forward}

Because the rules only alter the backward pass, the patched forward must be
unchanged. We verified max logit difference $= 0.0$ across all eight rule
combinations on the real model.

One implementation subtlety: the naive identity-rule form
\texttt{x * torch.sigmoid(x).detach()} differs from the fused \texttt{F.silu} by
roughly one ulp, which we measured at $2.3\times10^{-5}$ on final fp32 logits of the
4B model. We therefore implemented it as a custom autograd function whose forward
calls \texttt{F.silu} and whose backward returns $\sigma(x)$, giving an exactly
unchanged forward:

\begin{lstlisting}[language=Python]
class SiluWithSigmoidGrad(torch.autograd.Function):
    @staticmethod
    def forward(ctx, x):
        ctx.save_for_backward(x)
        return torch.nn.functional.silu(x)      # fused: forward is exact
    @staticmethod
    def backward(ctx, grad_output):
        (x,) = ctx.saved_tensors
        return grad_output * torch.sigmoid(x)   # the LRP identity rule
\end{lstlisting}

\subsection{The verification, and why it needs a noise floor}

Comparing our fitted lens to the released one requires knowing what agreement to
expect. An absolute tolerance is meaningless here, because two lenses fitted with
\emph{identical code} on \emph{different prompt draws} already disagree
substantially. So we measured that disagreement and used it as the bar.

\begin{center}
\begin{tabular}{@{}lrr@{}}
\toprule
\textbf{Comparison (4B, per-layer mean)} & \textbf{rel\_frob} & \textbf{corr} \\
\midrule
ours-R vs released-R                     & \textbf{0.0742} & 0.9961 \\
ours-J vs released-J                     & \textbf{0.1754} & 0.9737 \\
\emph{noise floor}: J-nf1 vs J-nf2       & \emph{0.3067}   & \emph{0.9405} \\
\emph{context}: released J vs released R & \emph{0.4470}   & \emph{0.8855} \\
\bottomrule
\end{tabular}
\end{center}

\noindent
where $\mathrm{rel\_frob} = \|A - B\|_F / \|B\|_F$, and the appended identity layer is
excluded.

Both of our lenses land \emph{under} the noise floor: our reimplementation differs
from theirs by less than two runs of the same code differ from each other. The last
row is the control in the other direction --- released J vs released R sit at 0.4470,
confirming the metric has the resolution to see a real difference when one exists.
Without it, small numbers everywhere could simply mean the metric is blind.

Behaviourally, on the full 496-item battery: ours-J $0.080$ vs released-J $0.082$;
ours-R $0.096$ vs released-R $0.098$ --- agreement to $0.002$.

\medskip
\noindent\fbox{\begin{minipage}{0.97\textwidth}
\textbf{Verdict: PASS.} The rebuilt fitting harness reproduces the released
artifacts. Everything downstream is licensed by this.
\end{minipage}}

% ===========================================================================
\section{Experimental design}

\subsection{The factorial}

Eight arms, one per subset of the three rules:

\begin{center}
\begin{tabular}{@{}lccc@{}}
\toprule
\textbf{Arm} & \cfg{ln} & \cfg{identity} & \cfg{half} \\
\midrule
\cfg{j} (baseline) & --- & --- & --- \\
\cfg{ln} & \checkmark & --- & --- \\
\cfg{identity} & --- & \checkmark & --- \\
\cfg{half} & --- & --- & \checkmark \\
\cfg{ln+identity} & \checkmark & \checkmark & --- \\
\cfg{ln+half} & \checkmark & --- & \checkmark \\
\cfg{identity+half} & --- & \checkmark & \checkmark \\
\cfg{r} (full \rlens) & \checkmark & \checkmark & \checkmark \\
\bottomrule
\end{tabular}
\end{center}

The endpoints \cfg{j} and \cfg{r} are the \jlens{} and \rlens{} respectively, which
gives the design two built-in sanity checks: \cfg{j} must reproduce our verified
\jlens, and \cfg{r} must land closest to the released \rlens.

\subsection{Controls}

\begin{itemize}[leftmargin=1.4em]
  \item \textbf{Identical prompts.} Every arm uses the same pile-10k rows, the same
        \texttt{t\_max}, the same \texttt{skip\_first}. The rules change the backward
        pass, not the data.
  \item \textbf{Identical recipe.} Same target layer, same estimator, same dtype.
  \item \textbf{Provenance verification.} Each fitted lens must carry the
        \texttt{config\_json} of the arm it is named for. A mislabelled sweep is the
        cheapest way to obtain a confidently wrong answer, so this is checked rather
        than assumed. All arms verified: correct \texttt{config\_json}, one shared
        fit commit, one shared prompt set.
  \item \textbf{Ordering for graceful degradation.} Single rules first
        ($\text{j} \to \text{ln} \to \text{identity} \to \text{half}$), then pairs,
        then the full lens. The single rules answer the headline question, so a
        truncated sweep still yields a coherent result.
\end{itemize}

\subsection{Two independent readouts}

\begin{enumerate}[leftmargin=1.4em]
  \item \textbf{Behavioural} --- pass@10 on the five official eval sets (multihop,
        multilingual, association, typo, poetry), per layer, with the correctness
        filter. This is the metric the original work reports.
  \item \textbf{Weight-space} --- per-layer cosine of each arm's $J_\ell$ against the
        released \rlens. Requires no eval at all.
\end{enumerate}

\noindent
We originally treated agreement between these as what makes an attribution credible.
Section~\ref{sec:discordance} reports that this expectation was wrong, which is
itself one of the findings.

% ===========================================================================
\section{Execution: what the runs actually cost}

\subsection{Calibration (measured, not estimated)}

Fitting cost scales with $d_{\text{model}} \times \text{params}$. We measured rather
than extrapolated:

\begin{center}
\begin{tabular}{@{}lll@{}}
\toprule
\textbf{Setting} & \textbf{Result} & \textbf{Note} \\
\midrule
\cfg{dim\_batch=8}  & 647 s/prompt, 99{,}209 MiB & measured on H200 143\,GB \\
\cfg{dim\_batch=16} & not attempted & would need $\sim$90\,GB activations; exceeds card \\
\cfg{dim\_batch=32} & not attempted & $\sim4\times$ the \cfg{dim\_batch=8} footprint \\
\bottomrule
\end{tabular}
\end{center}

\noindent
The $99$\,GB at \cfg{dim\_batch=8} is $54$\,GB of model plus $\sim45$\,GB of
activations, leaving no room to double the batch. This forced the 27B prompt count
down to $n=4$: $8\ \text{configs} \times 4\ \text{prompts} \times 647\,\text{s}
\approx 5.75$\,h.

\subsection{What completed}

\begin{center}
\begin{tabular}{@{}lccc@{}}
\toprule
\textbf{Model} & \textbf{Arms fitted} & \textbf{$n$} & \textbf{pass@10} \\
\midrule
Qwen3.5-4B  & 7 of 8 (missing \cfg{identity+half}) & 25 (= released recipe) & all arms \\
Qwen3.5-27B & 7 of 8 (missing \cfg{r})             & 4                      & all arms, 388 items \\
\bottomrule
\end{tabular}
\end{center}

\noindent
Neither model is individually complete, though between them all eight configurations
were fitted at least once. The 4B arm is the scientifically stronger one: $n=25$ is
\emph{exactly} the released recipe, against a baseline that provably reproduces the
published artifact.

% ===========================================================================
\section{Results}

\subsection{The effect being attributed}

Before asking which rule causes the \rlens{} advantage, we confirm the advantage
exists. Using only the \emph{released} lens pair --- independent of our fitting:

\begin{center}
\begin{tabular}{@{}lrrrr@{}}
\toprule
\textbf{27B, by eval set} & \textbf{logit} & \textbf{released J} & \textbf{released R} & \textbf{R $-$ J} \\
\midrule
typo         & 0.201 & 0.206 & 0.288 & $\mathbf{+0.0812}$ \\
multilingual & 0.085 & 0.136 & 0.165 & $\mathbf{+0.0284}$ \\
multihop     & 0.057 & 0.081 & 0.099 & $+0.0176$ \\
association  & 0.007 & 0.038 & 0.038 & $+0.0008$ \\
poetry       & 0.001 & 0.003 & 0.002 & $-0.0005$ \\
\midrule
\textbf{overall} & 0.0699 & 0.0929 & 0.1183 & $\mathbf{+0.0255 \pm 0.0027}$ \\
\bottomrule
\end{tabular}
\end{center}

\noindent
$+0.0255 \pm 0.0027$ SEM across 63 layers ($\approx 9\sigma$), rising to $+0.0362$
over the first half of the network. The effect is real and sizeable.

\textbf{Note which sets carry it.} \texttt{association} and \texttt{poetry} barely
move for \emph{any} lens, released ones included. See
Section~\ref{sec:floor}.

\subsection{Behavioural attribution}

\begin{center}
\begin{tabular}{@{}lrrrr@{}}
\toprule
\textbf{27B ($n=4$)} & \textbf{pass@10} & \textbf{$\Delta$ vs \cfg{j}} & \textbf{$t$(layers)} & \textbf{$t$(sets)} \\
\midrule
\cfg{j} (baseline) & 0.0943 & --- & --- & --- \\
\cfg{half}          & \textbf{0.1078} & $\mathbf{+0.0135}$ & $+4.5$ & $+1.8$ \\
\cfg{identity+half} & 0.1047 & $+0.0104$ & $+3.6$ & $+2.0$ \\
\cfg{ln+half}       & 0.1040 & $+0.0098$ & $+4.3$ & $+1.7$ \\
\cfg{ln+identity}   & 0.1017 & $+0.0074$ & $+2.7$ & $+2.1$ \\
\cfg{ln}            & 0.0879 & $-0.0064$ & $-3.0$ & $-1.8$ \\
\cfg{identity}      & 0.0878 & $-0.0065$ & $-5.2$ & $-1.6$ \\
\midrule
released-J & 0.0929 & & & \\
released-R & 0.1183 & & & \\
logit lens & 0.0699 & & & \\
\bottomrule
\end{tabular}
\end{center}

\begin{center}
\begin{tabular}{@{}lrrrr@{}}
\toprule
\textbf{4B ($n=25$)} & \textbf{pass@10} & \textbf{$\Delta$ vs \cfg{j}} & \textbf{$t$(layers)} & \textbf{$t$(sets)} \\
\midrule
\cfg{j} (baseline) & 0.1023 & --- & --- & --- \\
\cfg{half}        & \textbf{0.1219} & $\mathbf{+0.0196}$ & $+3.2$ & $+1.2$ \\
\cfg{r}           & 0.1202 & $+0.0179$ & $+4.0$ & $+1.5$ \\
\cfg{identity}    & 0.1075 & $+0.0052$ & $+2.3$ & $+1.4$ \\
\cfg{ln+identity} & 0.1026 & $+0.0003$ & --- & --- \\
\cfg{ln+half}     & 0.1021 & $-0.0002$ & --- & --- \\
\cfg{ln}          & 0.0754 & $-0.0269$ & $-3.7$ & $-1.4$ \\
\midrule
released-J & 0.1046 & & & \\
released-R & 0.1220 & & & \\
\bottomrule
\end{tabular}
\end{center}

\subsection{Normalised: fraction of the \rlens{} advantage recovered}

Dividing each rule's lift by the full released $\text{R}-\text{J}$ gap makes the two
models comparable despite different scales:

\begin{center}
\begin{tabular}{@{}lrr@{}}
\toprule
\textbf{Rule} & \textbf{27B} & \textbf{4B} \\
\midrule
\cfg{half}     & $\mathbf{53.0\%}$ & $\mathbf{111.9\%}$ \\
\cfg{ln}       & $-25.0\%$ & $-153.7\%$ \\
\cfg{identity} & $-25.6\%$ & $+29.9\%$ \\
\midrule
\emph{full advantage} & $+0.0254$ & $+0.0175$ \\
\bottomrule
\end{tabular}
\end{center}

\subsection{Headline}

\noindent\fbox{\begin{minipage}{0.97\textwidth}
\textbf{Of the three LRP rules, only the half-rule contributes to the \rlens's
advantage.} It recovers roughly half to all of the effect depending on model. The
LN-rule actively hurts on both. The identity-rule has no reliable effect --- it
changes sign between models ($+29.9\%$ at 4B, $-25.6\%$ at 27B).
\end{minipage}}

% ===========================================================================
\section{The weight-space discordance}
\label{sec:discordance}

We expected cosine-to-released-R to corroborate pass@10. It does not.

\begin{center}
\begin{tabular}{@{}lrrl@{}}
\toprule
\textbf{27B arm} & \textbf{$\Delta$cos to R} & \textbf{$\Delta$pass@10} & \textbf{agreement} \\
\midrule
\cfg{half}          & $+0.1853$ & $+0.0135$ & consistent \\
\cfg{identity+half} & $+0.1991$ & $+0.0104$ & consistent \\
\cfg{ln+half}       & $+0.1837$ & $+0.0098$ & consistent \\
\cfg{ln}            & $-0.1784$ & $-0.0064$ & consistent \\
\cfg{ln+identity}   & $-0.1445$ & $+0.0074$ & \textbf{sign wrong} \\
\cfg{identity}      & $+0.0333$ & $-0.0065$ & \textbf{sign wrong} \\
\bottomrule
\end{tabular}
\end{center}

\noindent
On 4B the failure appears elsewhere: \cfg{ln+half} is the \emph{second closest} arm
to the released \rlens{} in weight space ($\cos = 0.9858$, above \cfg{half}'s
$0.9811$) and yet scores $0.1021$ --- indistinguishable from the $0.1023$ baseline.

\subsection{Why this happens}

The mechanism is measurable. Two \jlens{}es fitted on \emph{disjoint} 25-prompt draws
differ by $\mathrm{rel\_frob} = 0.3067$, while an $n=4$ \jlens{} lands within
$0.0014$ pass@10 of an $n=25$ one. In other words:

\begin{center}
\emph{At these corpus sizes the lens is behaviourally saturated long before it is
geometrically converged.}
\end{center}

\noindent
Cosine-to-R measures a quantity still carrying heavy prompt-sampling noise; pass@10
measures one that stopped moving well before $n=25$. This also explains why the
verification in Section 2 had to be judged against a noise floor: an absolute
weight-match criterion would have been unmeetable at $n=25$ no matter how correct the
implementation.

\medskip
\noindent\fbox{\begin{minipage}{0.97\textwidth}
\textbf{Secondary finding.} Weight-space similarity to the released \rlens{} does not
predict lens quality, and gets the sign wrong for two arms. This is a contribution
beyond the original work, and unlike the attribution it is \emph{immune to the
corpus-size critique} --- both readings are taken on the same lenses, so lens quality
does not enter.
\end{minipage}}

\subsection{The $n=4$ noise floor, and which geometry readings survive it}
\label{sec:noisefloor}

The obvious objection to every weight-space number above is that each arm is a single
$n=4$ draw. We measured the alternative directly: a second \cfg{j} lens, identical
config and identical $n$, fitted on a \emph{disjoint} draw (pile-10k rows $[25,29)$
against the sweep's $[0,4)$). Its disagreement with the sweep's \cfg{j} is variation
attributable to prompt sampling alone, with no rule difference involved.

\begin{center}
\begin{tabular}{@{}lr@{}}
\toprule
\textbf{Two 27B \jlens{}es, $n=4$, disjoint draws} & \\
\midrule
$\mathrm{rel\_frob}$ mean (max)              & 0.786 (1.799) \\
$\mathrm{corr}$ mean (min)                   & 0.676 (0.152) \\
\textbf{noise floor on the cosine scale}     & $\mathbf{0.0574}$ \\
\bottomrule
\end{tabular}
\end{center}

\noindent
For context, the 4B floor at $n=25$ is $\mathrm{rel\_frob} = 0.307$. At $n=4$ it is
$0.786$, roughly $2.6\times$ worse --- while behaviour moved by only $0.0014$ pass@10
over the same reduction. That is the saturation asymmetry of
Section~\ref{sec:discordance} quantified at both ends.

Measuring each arm against the floor:

\begin{center}
\begin{tabular}{@{}lrrl@{}}
\toprule
\textbf{Arm} & \textbf{$\Delta$cos} & \textbf{$\times$ floor} & \textbf{verdict} \\
\midrule
\cfg{identity+half} & $+0.1991$ & $3.5$ & clears \\
\cfg{half}          & $+0.1853$ & $3.2$ & clears \\
\cfg{ln+half}       & $+0.1837$ & $3.2$ & clears \\
\cfg{ln}            & $-0.1784$ & $3.1$ & clears \\
\cfg{ln+identity}   & $-0.1445$ & $2.5$ & clears \\
\cfg{identity}      & $+0.0333$ & $0.6$ & \textbf{within noise} \\
\bottomrule
\end{tabular}
\end{center}

\noindent
Two things fall out, and both confirm conclusions we had previously reached on
judgement alone.

\begin{enumerate}[leftmargin=1.4em]
  \item \textbf{\cfg{identity} is the only arm that fails to clear the floor} --- and it
        is precisely the arm whose geometry mispredicted its behaviour in sign
        ($+0.0333$ cosine against $-0.0065$ pass@10). Its weight-space reading was
        never signal.
  \item \textbf{The three half-containing arms are mutually indistinguishable.}
        Pairwise gaps are $0.0016$--$0.0154$, i.e.\ $0.03$--$0.27\times$ the floor.
        They cannot be ranked in weight space, which is exactly the restriction we
        placed on them behaviourally.
\end{enumerate}

\noindent
So the geometry column is not worthless --- \cfg{half} and \cfg{ln} clear the floor by
more than $3\times$, and their signs match pass@10. It is unreliable specifically where
effects are small, which is where its two sign errors occurred. The
\cfg{ln+identity} case remains the genuine anomaly: it clears the floor at $2.5\times$
and is \emph{still} wrong in sign, which no amount of sampling noise explains and which
we cannot account for.

% ===========================================================================
\section{Caveats}

\subsection{Corpus size}
\label{sec:corpus}

The \jlens{} documentation states: \emph{``The paper's lenses use 1000 sequences of
128 tokens from a pretraining-like corpus. Quality saturates quickly (\S9.3);
$\sim$100 prompts is usable.''}

The released \rlens{} artifacts used \textbf{25}. This is below the paper's own
guidance, and is a property of the artifacts we are replicating rather than a defect
in the replication --- our 4B arm matches their $n$ exactly.

Our 27B arms used $n=4$, which is $6\times$ under the released recipe. We bounded its
cost directly rather than assuming:

\begin{center}
\begin{tabular}{@{}lll@{}}
\toprule
\textbf{Comparison} & \textbf{Isolates} & \textbf{Gap} \\
\midrule
our \cfg{j} ($n{=}4$) vs released-J ($n{=}25$), 27B & implementation $+$ $n$ drop & $+0.0014$ \\
our \cfg{j} ($n{=}25$) vs released-J ($n{=}25$), 4B & implementation alone & $-0.0023$ \\
\bottomrule
\end{tabular}
\end{center}

\noindent
Dropping from 25 prompts to 4 moved pass@10 by \emph{less} than our implementation's
own reproduction error at matched $n$. Both are $\approx 0.002$, roughly $10\times$
smaller than the half-rule's effect. \textbf{Corpus size is not what threatens the
headline} --- but it does undercut every weight-space number, consistent with
Section~\ref{sec:discordance}.

\subsection{Floor effects in the eval battery}
\label{sec:floor}

On both models, \texttt{association} and \texttt{poetry} sit near zero for
\emph{every} lens including the released ones. At 27B their released
$\text{R}-\text{J}$ gaps are $+0.0008$ and $-0.0005$. They carry no signal.

Consequently every per-rule lift reported here is an attribution of the
\emph{typo/multilingual} effect, and should be quoted that way rather than as a claim
about the five-set battery. This applies equally to the original work's claims.

\subsection{Which uncertainty is being reported}

The same numbers give very different confidence depending on the sampling unit:

\begin{center}
\begin{tabular}{@{}lrr@{}}
\toprule
\textbf{4B rule} & \textbf{$t$ across layers} & \textbf{$t$ across eval sets} \\
\midrule
\cfg{half}     & $+3.2$ & $+1.2$ \\
\cfg{identity} & $+2.3$ & $+1.4$ \\
\cfg{ln}       & $-3.7$ & $-1.4$ \\
\bottomrule
\end{tabular}
\end{center}

\noindent
Across-layer SEM asks \emph{``is the effect consistent down the network?''} (yes,
strongly). Across-set SEM asks \emph{``does it generalise across tasks?''} (not
established). Both are honest; reporting only the first overstates the result. Note
also that layers are \emph{not} independent draws --- a bad prompt sample shifts every
layer together --- so the across-layer $t$ provides no protection against sampling
error in the fitting corpus.

\subsection{Item-level uncertainty is unavailable}

The plan called for bootstrap CIs over items. The per-layer CSV stores aggregated
fractions, discarding per-item outcomes, so item-level CIs cannot be computed from
the saved artifacts. Retaining per-item outcomes would be a cheap change and is the
single highest-value fix for a rerun.

\subsection{Incomplete arms}

4B is missing \cfg{identity+half}; 27B is missing \cfg{r}. The 27B omission
specifically disables the design's sanity check that \cfg{r} lands closest to the
released \rlens.

\subsection{What cannot be claimed}

\begin{itemize}[leftmargin=1.4em]
  \item Any ranking among \cfg{half} / \cfg{identity+half} / \cfg{ln+half} at 27B ---
        gaps of $\sim0.003$, only $\sim2\times$ the measured reproduction error.
  \item That any 27B lens of ours approximates a properly fitted lens.
  \item Generalisation across the full eval battery (across-set $t$ is $1.2$--$2.1$).
  \item Anything about specific weight values.
\end{itemize}

% ===========================================================================
\section{Things we got wrong}

Recorded because several were near-misses that would have produced confident wrong
answers rather than loud failures.

\subsection{Silent wrong answers}

\begin{description}[leftmargin=1.4em,style=nextline]
\item[Directory naming would have emptied the attribution table.]
  \texttt{rlens fit} writes the two endpoints to \texttt{j-lens/} and \texttt{r-lens/}
  but every other arm to its bare config name. Code that assumed one convention
  silently resolved nothing for half the arms.

\item[The fix for it was worse than the bug.] Our first correction resolved the output
  path \emph{before} the fit ran --- when neither directory exists yet, so the fallback
  fired for every arm. Path resolution has to happen after the fit.

\item[Pandas drops MultiIndex names on CSV round-trip.] \texttt{to\_csv} does not
  preserve MultiIndex \emph{names}, so every later
  \texttt{groupby(level="lens")} failed silently. Caught by validating the analysis
  against \emph{synthetic} lenses with a planted signal, before any real results
  existed --- a practice worth keeping.

\item[The noise-floor script would have reported ``zero noise''.] Our first version
  copied the fitted lens from \texttt{j-lens/}, not knowing that
  \texttt{-{}-draw nf1} writes to \texttt{j-lens-nf1/}. It would have banked a duplicate
  of the \emph{primary} lens under the noise-floor name and reported a floor of
  $\approx 1.000$ --- ``no sampling noise at all,'' entirely wrong and entirely
  plausible-looking.

\item[A verdict threshold that flattered small effects.] Our report generator scored
  each rule against a fraction of the maximum, which labels a $+0.03$ rule a
  ``carrier'' whenever the true carrier happens to be large. Replaced with a dominance
  test: top rule, positive, at least twice the runner-up.

\item[Hardcoded caveats going stale.] The figure page asserted ``27B is weight-space
  only'' as fixed prose. That sentence became false the moment the eval was recovered,
  while the panels silently changed underneath it. Caveats are now derived from what
  is on disk.
\end{description}

\subsection{Infrastructure}

\begin{description}[leftmargin=1.4em,style=nextline]
\item[\texttt{OSError: [Errno 122] Disk quota exceeded} --- the root cause of everything.]
  \texttt{df} reported 299\,TB free while writes failed. A per-volume quota was
  exhausted by 53\,GB of \texttt{fit.ckpt.pt} checkpoint files (one per arm, $\sim$7\,GB
  each at 27B). This single fault explained a set of unrelated-looking symptoms: fits
  dying silently, \texttt{scp} reporting \texttt{close remote: Failure}, and the eval
  computing for an hour then producing no file. Deleting the checkpoints restored
  writes.

\item[Two monitoring failures, both costing hours.] The first watched a \emph{log file},
  where ``stalled'' is indistinguishable from ``slow.'' The second matched the
  \emph{wrapper} shell rather than the actual worker process, so a dead job looked
  alive. Monitors must watch the worker process and must emit on every terminal
  state, not only the success signature.

\item[\texttt{run\_passk} printed nothing at all.] A 27B eval over ten lenses runs for
  an hour with no output, which is what made the above failures invisible. Now emits
  one flushed line per scored item.

\item[\texttt{pkill -f chain\_r\_arm} killed its own SSH shell,] because the pattern
  matched the command line containing it. Everything after the \texttt{pkill} silently
  did not run, leaving no job armed to bank results.

\item[\texttt{git checkout -- .} destroyed 2.3\,h of results.] Replaced with a sync
  script that resets code paths only.

\item[A hardcoded model guard.] \texttt{cmd\_fit} contained
  \texttt{if args.model != "qwen3.5-4b": raise SystemExit}, which killed the first 27B
  calibration. The wrapper script misreported this as an OOM.
\end{description}

\subsection{What we would do differently}

\begin{enumerate}[leftmargin=1.4em]
  \item Retain per-item eval outcomes, enabling item-level bootstrap CIs.
  \item Report across-set uncertainty alongside across-layer from the start.
  \item Fit the noise-floor arm \emph{first}, not last --- it calibrates every other
        number and costs one arm's compute.
  \item Delete checkpoints on arm completion; they are resume state, not artifacts.
  \item Drop \texttt{association} and \texttt{poetry}, or replace them with sets that
        are not at floor, and spend the compute on more prompts instead.
\end{enumerate}

% ===========================================================================
\section{Reproduction}

\begin{lstlisting}[language=bash]
# fit one arm (any subset of the three rules)
uv run rlens fit --model qwen3.5-27b --lens half --n 4 --dim-batch 8

# the full ordered sweep, with per-arm copy-off
./scripts/run_lrp_sweep.sh qwen3.5-27b 4 8

# pass@10 over every fitted arm + released + logit lens
uv run rlens eval --model qwen3.5-27b

# analysis (CPU only, from the committed CSVs -- no GPU, no lens weights)
uv run python scripts/report_lrp.py
uv run python scripts/plot_lrp.py

# the noise floor: same config, same n, disjoint prompt draw
uv run rlens fit --model qwen3.5-27b --lens j --draw nf1 --n 4 --dim-batch 8
\end{lstlisting}

\paragraph{Artifacts.}
\texttt{results/lrp/}: \texttt{lrp\_qwen3.5-\{4b,27b\}.md} (reports),
\texttt{figures\_lrp.html} (7 panels),
\texttt{passk\_per\_layer\_*.csv} (per-layer pass@10),
\texttt{lrp\_geometry\_*.csv} (per-layer cosines),
\texttt{lrp\_labels\_*.csv} (provenance verification),
\texttt{lrp\_timing\_qwen3.5-27b.json} (measured calibration).

\end{document}
